クリフディメンションの背面飛び移りを，3\,cm の突起の有限な把持容量と前腕の干渉を含む力学課題として定式化し，
\gridN 条件の多相軌道最適化から必要把持容量を求めた．必要容量は体重の \capBWmin〜\capBWmax 倍で，装置周期に依らず，捕捉後の保持が決める．
さらに，最適化の双対変数（価値の勾配と制約の価格）を教師にする双対場学習を提案し，同じ物理の環境で検証した．
双対ラベルは勾配の精度，候補の順位付け，離手率を改善したが，閉ループの成功率は \clSmpcSuccInt\,\%のままであった．
参照解の再生も成功率 \clReplaySuccInt\,\%にとどまり，難しさは制約の境界上を通る最適解の実行にある．
